\section{Regularisation and Preconditioning}\label{sec:conditioning}

Both the inner CG rate (Proposition~\ref{prop:cg}) and the projected-gradient rate
(Proposition~\ref{prop:proj}) are governed by $\kappa = \kappa(A)$, so anything that
compresses the spectrum of $A$ shortens the solve. Two mechanisms are available, and they
coincide in a convenient way.

\paragraph{Regularising split.}
Replace $A$ by the convex combination
\begin{equation}\label{eq:regsplit}
  A_\alpha = (1-\alpha)\,A + \alpha\,R^\top R, \qquad \alpha \in (0,1),\; R^\top R \succ 0,
\end{equation}
a ridge/Tikhonov regularisation of $A$ toward the SPD target $R^\top R$. When $A = M^\top M$
this is again a Gram matrix, of the stacked operator
\[
  \tilde{M} = \begin{pmatrix} \sqrt{1-\alpha}\,M \\ \sqrt{\alpha}\,R \end{pmatrix},
  \qquad \tilde{M}^\top\tilde{M} = A_\alpha,
\]
so it is simply the Gram backend of $\tilde M$ (the regularised operator of
Section~\ref{sec:operators}) and CG runs on $\tilde{M}_{\mathcal{F}}$ with no change. The
augmented-matrix device is standard in least-squares
regularisation~\cite{golub2013,lawson1974}.

\begin{proposition}[Condition number under the regularising split]\label{prop:regcond}
  For $A_\alpha$ as in~\eqref{eq:regsplit}, Weyl's inequality~\cite{hornjohnson2013} gives
  \[
    \kappa(A_\alpha) \;\leq\;
    \frac{(1-\alpha)\lambda_{\max}(A) + \alpha\,\lambda_{\max}(R^\top R)}
         {\alpha\,\lambda_{\min}(R^\top R)}.
  \]
  The bound is strictly decreasing in $\alpha$; as $\alpha \to 1$ it approaches
  $\kappa(R^\top R)$, which equals $1$ for $R = \sqrt{c}\,I$. For the scaled-identity target
  $R^\top R = \bar\lambda I$ it simplifies to
  $\kappa(A_\alpha) \le \tfrac{1-\alpha}{\alpha}\,\tfrac{\lambda_{\max}(A)}{\bar\lambda} + 1
  = O\!\big((1-\alpha)/\alpha\big)$.
\end{proposition}

The split does double duty. It lowers $\kappa$ and hence the iteration count of both solvers
(Propositions~\ref{prop:cg} and~\ref{prop:proj}), and, because $A_\alpha \succ 0$ strictly
whenever $\alpha > 0$, even when $A$ is only positive semidefinite (a rank-deficient
least-squares operator with $m < n$), it restores the $P$-matrix property on which
Theorem~\ref{thm:termination} depends. A strictly positive $\alpha$ is therefore the natural
setting for the active-set loop: it simultaneously conditions the inner solve and guarantees
that every principal submatrix is invertible.

\paragraph{Preconditioning.}
When the operator must be solved at its given $A$, a symmetric positive definite
preconditioner $P \approx A$ replaces the CG rate's $\kappa(A)$ by $\kappa(P^{-1}A)$; standard
choices (Jacobi, incomplete Cholesky, or a low-rank-plus-diagonal approximation built from a
few dominant eigenpairs) apply unchanged inside the active-set loop, since preconditioned CG
still accesses $A_{\mathcal{F}}$ only through its mat-vec~\cite{golub2013,benzi2005}. The
regularising split~\eqref{eq:regsplit} may itself be read as an implicit preconditioner: it
changes the problem, but toward a target the application deems statistically or structurally
preferable, so the conditioning gain comes at no modelling cost. The companion
papers~\cite{schmelzer2026matrixfree,schmelzer2026rmt} develop concrete instances (a
scaled-identity split and a low-rank-plus-identity target inverted by the Woodbury identity)
and quantify the resulting iteration savings on their structured operators.
